\documentclass[11pt]{article}
\usepackage[a4paper,margin=19mm]{geometry}
\usepackage{fontspec}
\setmainfont{lmroman10-regular.otf}
\usepackage{amsmath,amssymb,amsthm,mathtools,amscd,xcolor,graphicx,hyperref}
\hypersetup{hidelinks}
\numberwithin{equation}{section}
\newtheorem{theorem}{Theorem}[section]
\newtheorem{lemma}[theorem]{Lemma}
\newtheorem{proposition}[theorem]{Proposition}
\newtheorem{corollary}[theorem]{Corollary}
\newtheorem{claim}[theorem]{Claim}
\newtheorem{fact}[theorem]{Fact}
\newtheorem{conjecture}[theorem]{Conjecture}
\theoremstyle{definition}
\newtheorem{definition}[theorem]{Definition}
\newtheorem{example}[theorem]{Example}
\newtheorem{remark}[theorem]{Remark}
\newtheorem{question}[theorem]{Question}
\theoremstyle{plain}
\newcommand{\C}{\mathbb C}
\newcommand{\R}{\mathbb R}
\newcommand{\Q}{\mathbb Q}
\newcommand{\Z}{\mathbb Z}
\newcommand{\N}{\mathbb N}
\newcommand{\eee}{\mathrm e}
\newcommand{\cclass}[1]{\textsf{#1}}
\newcommand{\cproblem}[1]{\textsc{#1}}
\newcommand{\ie}{i.\,e.}
\newcommand{\Ie}{I.\,e.}
\newcommand{\eg}{e.\,g.}
\newcommand{\Eg}{E.\,g.}
\newcommand{\phd}{Ph.\,D.}
\newcommand{\msc}{M.\,S.}
\newcommand{\bsc}{B.\,S.}
\newcommand{\expref}[2]{\hyperref[#2]{#1~\ref{#2}}}
\newcommand{\expeqref}[2]{\hyperref[#2]{#1~\eqref{#2}}}
\newcommand{\exprefsilent}[2]{\hyperref[#2]{#1}}
\providecommand{\etal}{et~al.}
\usepackage{eprint}
\usepackage[linesnumbered,boxed,ruled,vlined]{algorithm2e}
\usepackage{thm-restate}
\usepackage{commath}


\DeclareMathOperator*{\E}{\mathbb{E}}
\newcommand{\Braket}[1]{\langle#1\rangle}
\newcommand{\poly}{\text{poly}}
\newcommand{\W}{\mathbf{W}}
\newcommand{\I}{\mathbf{I}}
%
\newcommand{\cC}{\mathcal{C}}
\newcommand{\cF}{\mathcal{F}}
\newcommand{\floor}[1]{\lfloor#1\rfloor}

%
\newcommand{\warn}[1]{#1}
\title{Slice-Fourier distinguishers for negation-limited weak pseudorandom functions}
\author{Zhihuai Chen, Siyao Guo, Qian Li, Chengyu Lin, Xiaoming Sun}
\date{doi:10.4086/toc.2024.v020a002; CC BY 3.0}
\begin{document}
\maketitle
\noindent\textit{Editorial note.} Original statement, author proof, algorithms and bibliography retained; standalone wrapper only.\par\medskip
\noindent\textit{Allowed standard prerequisites (editorial).}\begin{itemize}
\item Boolean functions, probability and basic Fourier analysis. Original source slice basis, k-monotonicity and cited concentration machinery retained.. No collector proof or repaired assumption.
\end{itemize}\medskip
\section{Introduction}

One significant goal in the area of cryptography is to understand how simple cryptography can be.  This motivates the study of low complexity cryptography which explores the possibility of implementing cryptographic primitives in low complexity classes. This line of research inherently lies at the intersection of computational complexity and cryptography. It links core problems in both areas and has become an essential source of new perspectives for both areas.

In this work, we continue this line of research and focus on \emph{pseudorandom functions} (PRFs) in \emph{negation-limited computation}. We start by introducing pseudorandom functions and negation-limited computation before connecting them to explain the main motivation of our work.

\vspace{2ex}
\noindent\textbf{Pseudorandom functions.} Pseudorandom functions (PRFs)~\cite{GGM} are fundamental primitives in symmetric cryptography. In particular, they yield direct solutions to most central goals of symmetric cryptography, such as encryption, authentication and identification. They are well studied in the theoretical community and widely used in practice.

As lightweight (computationally limited) devices become popular, the efficiency of cryptographic implementations
also %
becomes %
increasingly significant.  To obtain a better tradeoff between efficiency and security, a weaker notion of PRFs called \emph{weak pseudorandom functions} (See \expref{Definition}{def:wprfs}) has been considered. A distinguisher for a 
family of %
weak PRFs aims to distinguish a random member of the family from a truly random function after observing a number of random samples $(x_1,f_s(x_1)),\dots, (x_m,f_s(x_m))$ where $x_1,\dots,x_m$ are independent uniformly random strings from $\{0,1\}^n$ and $f_s:\{0,1\}^n\rightarrow\{0,1\}$ is the function in question. Weak PRFs suffice for many key applications such as encryption and authentication in symmetric cryptography. More importantly, weak PRFs may allow for significant gains in efficiency.  Akavia et al.~\cite{ABGKM14} pointed out weak PRFs have the potential to bypass the limitations of PRFs in low depth circuits. In particular, they provided candidate weak PRFs in a class of low depth circuits where PRFs provably cannot exist. This raises the following natural questions.
\begin{quote} {\it
Can weak PRFs bypass the limitations of PRFs in other low complexity classes? }
\end{quote}

Besides cryptography, another important motivation for the study of low complexity PRFs comes from explaining the difficulties of obtaining circuit lower bounds and learning algorithms. We refer interested readers to the survey by Bogdanov and Rosen~\cite{BR17}. 


\vspace{2ex}
%
\noindent\textbf{Negation-limited computation.} %
The power of negations is a mystery in complexity theory. One of the main difficulties in proving lower bounds on circuit size using AND, OR, NOT gates is the presence of negation gates: the best such lower bound is linear, whereas if no negation gates are allowed, exponential lower bounds are known~\cite{razborov1985lower,andreev1985method,alon1987monotone,tardos1988gap,berg1999symmetric,harnik2000higher}. In 1958, Markov~\cite{Mark} observed that every Boolean (even multiple-output) function of $n$ variables can be computed by a circuit with only $\log{n}$ negation gates.  In other words, the 
%
%
exponential gap between monotone computation and non-monotone computation exists due to as few as $\log{n}$ negations.

Besides circuit complexity, the divide between monotone and non-monotone computation exists in general:  while we usually have a fairly good understanding of the monotone case, many things may fail to hold when negation gates are allowed. Aiming at bridging the gap between monotone and non-monotone computation, a body of recent work studies negation-limited computation from multiple angles including learning~\cite{BCOST15}, cryptography~\cite{GMOR15}, Boolean formulas~\cite{GK15,Ros15}, property testing~\cite{CGGKW17,GKW19}, Boolean function conjectures~\cite{LZ17}. Although the above works extend many results in monotone cases to as many as $O(\log n)$ negations, they also leave open several surprisingly basic questions about a single negation ranging from weak learning algorithms to the structure of their Fourier spectrum. More surprisingly, in the context of property testing, a single negation can be exponentially harder than the monotone case~\cite{CGGKW17,GKW19}.  Our understanding of a single negation remains largely a mystery.

When the circuit size is not of interest, the classes of circuits with $\log k$ negations are captured by the class of so-called $k$-monotone functions where each function in the family can be written as the parity of $k$ monotone
%
functions (see  %
\expref{Section}{sec:k-mon}). To simplify the presentation, we will use $k$-monotone functions instead of circuits with $\log{k}$ negations in some of our discussions.

\vspace{2ex}
%
\noindent\textbf{PRFs in negation-limited computation.} %
Can pseudorandom functions be computed by a few negations? For pseudorandom functions, we have a fairly good understanding. Guo et al.~\cite{GMOR15} showed that PRFs are inherently highly non-monotone and require $\log{n}-O(1)$ negations, which is optimal up to an additive constant. However, the answer to weak PRFs is unsatisfying. Guo et al.~\cite{GMOR15} observed that weak PRFs cannot be monotone due to the weak learner for monotone functions by Blum et al.~\cite{BABL98}. For general $k$, the best distinguisher, due to Blais et al.~\cite{BCOST15}, runs in time $n^{O(k\sqrt{n})}$.  Therefore even for a single negation (\ie, $k=2$), the best distinguisher runs in time $n^{O(\sqrt{n})}$.

The above results demonstrate two strong separations. In negation-limited computation,
weak PRFs have the potential to be much simpler than PRFs: 
for all we know, %
even a single negation
may have %
%
%
%
%
$n^{\Omega(\sqrt{n})}$ %
hardness whereas PRFs cannot exist. From the angle of weak PRFs, the hardness gap between even a single negation and monotone 
%
could %
be %
as large as 
%
$n^{\Omega(\sqrt{n})}$.  %
%
These separations are our main motivation to connect them together to study negation-limited weak PRFs. 

\subsection{Our results}
Before presenting our main results, we define \emph{weak pseudorandom functions} and \emph{weak learning under uniform distribution}.
%
%
\begin{definition}[Weak pseudorandom functions]\label{def:wprfs}
Let $S$ be a distribution over $\{0,1\}^m$ and $\{F_s:\{0,1\}^n\rightarrow\{0,1\}\}$ be a family of functions indexed by string $s$ in the support of $S$. We say
that  %
$\{F_s\}$ is a
family of %
$(c,\epsilon)$-secure weak pseudorandom functions (wPRFs) if for every
%
%
%
Boolean circuit $D$  %
of size at most $c$,
\begin{align}\label{ineq:def1}
\abs{\Pr_s[D^{F_s}\mbox{ accepts}]-\Pr_R[D^R\mbox{ accepts}]}\leq \varepsilon,
\end{align}
where $s$ is distributed according to $S$, $R$ is a function sampled uniformly at random from the set of all functions from $\{0,1\}^n$ to $\{0,1\}$, and $D^h$ denotes the execution of $D$ with random oracle access to a Boolean function $h:\{0,1\}^n\rightarrow\{0,1\}$. In other words, the distinguisher $D^h$ only has access to random examples of the form $(x,h(x))$ where $x$ is uniformly distributed over $\{0,1\}^n$. \warn{The two probabilities in \expeqref{Equation}{ineq:def1} are both also over the
%
random samples $x$.}  %
\end{definition}

%
%
%
%
%
%
%
\begin{definition}[Weak learning under the uniform distribution] %
We say that an algorithm $A$ weakly learns a family
%
$\cF$ of Boolean functions %
under the uniform distribution, %
%
%
%
if given any $f\in \cF$ it accesses pairs $(x,f(x))$ where $x$ is picked from $\{0,1\}^n$ uniformly at random, and then %
%
%
%
%
%
outputs a hypothesis $h$ such that with high probability \warn{(over the random samples and the randomness of $A$)}
%
\begin{equation} %
 \Pr_{x \sim U_{\{0, 1\}^n}}\left[ f(x) \neq h(x) \right] \le \frac{1}{2}-\frac{1}{\poly(n)}   
%
\end{equation} %
where $U_{\{0,1\}^n}$ is the uniform distribution over $\{0,1\}^n$.
\end{definition}

A weak learner works slightly better than random guessing. But from this small
advantage, if %
it is  %
non-negligible,
one can naturally derive an efficient distinguisher against
a  %
random function.
Any weak learner explicitly gives an attack on  %
%
candidate families of weak pseudorandom functions.  %
Conversely,  %
families of %
%
%
weak pseudorandom functions are hard to learn.
Our main result is new distinguishers for negation-limited %
families of %
weak pseudorandom functions. Our results hold for inefficient circuits and are stated in terms of $k$-monotone functions.


\begin{restatable}{theorem}{distinguisher}\label{thm:distinguisher}
Any family of $k$-monotone functions can be distinguished from uniformly random functions in $\exp\left(O(n^{1/3}(k\log n)^{2/3})\right)$ time.
In other words, any family of $k$-monotone functions is not a
	  $\left(\exp\left(O(n^{1/3}(k\log n)^{2/3})\right),1/3\right)$-secure
	  weak pseudorandom family.
\end{restatable}

The previous best distinguisher for $k$-monotone weak PRFs is the learning algorithm by Blais et al.~\cite{BCOST15} which runs in $\exp\big(O(n^{1/2}k\log n)\big)$ time. Our result improves
this bound by %
an $\Omega\big(n^{1/6}(k\log n)^{1/3}\big)$ factor in the exponent.

\expref{Theorem}{thm:distinguisher} implies that
%
%
exponentially secure weak PRFs require  %
%
$\log{n}-O(\log\log{n})$ negations, which is optimal up to an additive $O(\log\log{n})$ term. Therefore, weak PRFs cannot bypass the limitations of PRFs in terms of achieving exponential security.

%
\expref{Theorem}{thm:distinguisher} \ also \ implies \ that \
%
$1$-negation \ functions \ %
can \ be \ distinguished \ in \
$\exp\big(O(n^{1/3}\log^{2/3}{n})\big)$ time. Therefore,  unlike testing $1$ negation (using $1$-sided non-adaptive tester)~\cite{GKW19} and learning $1$ negation to high accuracy~\cite{BCOST15}, distinguishing $1$ negation
%
does not %
suffer from the $\exp(\sqrt{n})$ barrier.

It is natural to ask if we can leverage the distinguisher to a learning algorithm. Our second result gives weak learning algorithms for $k$-monotone functions under the uniform distribution.
\begin{restatable}{theorem}{learner}\label{thm:learner}
	$k$-monotone functions are weakly learnable in time $\exp\left(O\big(k\sqrt{n\log n}\big)\right)$.
\end{restatable}

Our result slightly improves the previous best weak learner due to Blais et al.~\cite{BCOST15}, by a $\Omega\big(\sqrt{\log n}\big)$ factor in the exponent.

\warn{We conjecture that both \expref{Theorem}{thm:distinguisher} and \expref{Theorem}{thm:learner} are not tight. However, we believe that any further improvement of our results, even for a single negation, 
%
requires %
completely new techniques or proving rather hard conjectures which seem out of reach. See \expref{Section}{sec:discussion} for more details.}

\vspace{2ex}
%
\noindent\textbf{Our techniques.} %
Blais et al.~\cite{BCOST15} showed a Fourier concentration of $k$-monotone functions on low degree monomials, by bounding the total
influence of $k$-monotone functions. Then they apply the
``Low-Degree Algorithm'' established by
Linial, Mansour, and Nisan~\cite{LMN93} to learn $k$-monotone functions. One natural idea to improve their learning algorithm is to show Fourier concentration on lower levels. However,
their influence bound is tight and even for monotone functions, we cannot show  concentration bound on fewer than $\Omega(\sqrt{n})$ levels ~\cite{DFTWW14}, which will require at least $n^{\Omega(\sqrt{n})}$ time by applying the ``Low-Degree Algorithm''.

Our main technique is using Fourier analysis on slices~\cite{Filmus,srinivasan2011symmetric}. Although the Fourier concentration on the Boolean cube cannot be improved, we show some ``middle'' slices of $k$-monotone functions can have much stronger Fourier concentration.  Then by adapting the ``Low-Degree Algorithm'' to the slices, we obtain a distinguisher with significantly improved running time. Our weak learner is a simple variant of the ``Low-Degree Algorithm'' on slices.

Fourier analysis on slices has a notion of total influence which allows us to show Fourier concentration on a slice in a similar way. We give an upper bound on the sum of total influences for all ``middle'' slices of any $k$-monotone function. It implies the existence of a ``middle'' slice function with small total influence, and therefore good concentration. Then we optimize the number of ``middle'' slices to be analyzed to get an efficient algorithm.

\vspace{2ex}
%
%
\noindent\textbf{Organization of the paper.} %
We begin with %
%
basic notation  %
in Section 2, then present the structural results for $k$-monotone functions
in Section 3. In Sections 4 and 5, we present the distinguisher and
the  %
%
%
weak learner, respectively.  %

\section{Preliminaries}
%
%
%
\subsection{Notation, terminology: sets, strings, slices}
\warn{In this paper, all the logarithms are base 2.}

As usual, for $n\in \N$ we write $[n]=\{1,2,\dots, n\}$.  The set of
$k$-subsets (subsets of size $k$) of $[n]$ is denoted $\binom{[n]}{k}$.

We refer to the set $\{0,1\}^n$ as the $n$-dimensional \emph{Boolean hypercube};
its elements are the \emph{$(0,1)$-strings} of length $n$.   
For $0\le r\le n$, the \emph{$r$-slice} of the $n$-cube is the set
\begin{equation}
\{(x_1,\ldots,x_n)\in\{0,1\}^n\,:\,\sum_i{x_i}=r\}\,.
\end{equation}
    
We identify subsets $A\subseteq [n]$ with their indicator strings
$x_A=(x_1,\dots,x_n)$ defined by
\begin{equation}
   x_i = \begin{cases} 1 & \mbox{if } i\in A\\
                       0 & \mbox{if } i\notin A\,.
         \end{cases}
\end{equation}
With this identification, the $r$-slice of the $n$-dimensional
Boolean hypercube becomes identical with the set $\binom{[n]}{r}$.   
%

%
\subsection{Alternating number, negation complexity, $k$-monotone functions}
\label{sec:k-mon}

%
%
For any two strings $x, y \in \{0, 1\}^n$, we say $x \prec y$
(or $y \succ x$) %
%
if   %
$x \neq y$ and $x_i \le y_i$ for all $i \in [n]$. 
A chain $X = (x^1, x^2, \dots, x^\ell)$ of length $\ell$ is an increasing sequence of strings in $\{0, 1\}^n$
where $x^i \prec x^{i+1}$ for $i \in [\ell - 1]$.
For a Boolean function $f:\{0,1\}^n\rightarrow \{0, 1\}$, we define the 
\emph{alternating number} %
of $f$
on chain $X$ to be the number of value flips on this chain:
\begin{equation}
a(f, X) = \left| \{i \in [\ell - 1] : f(x^i) \neq f(x^{i+1})\}\right|.
\end{equation}
Let $\cC$ be the set of all chains on $\{0, 1\}^n$, the alternating number of $f$ is
\begin{equation}
a(f) = \max_{X \in \cC} a(f, X).
\end{equation}
Note that the alternating number of a monotone function is no more than 1.

A celebrated result of Markov connects the alternating number of a Boolean function $f$ to the
negation complexity $\mathbf{N}(f)$ -- the minimum number of negation gates required in any
Boolean $\sf{AND}-\sf{OR}$ circuits to compute $f$.

\begin{theorem}[Markov's Theorem~\cite{Mark}]\label{thm:markov}
Let $f:\{0,1\}^n \rightarrow \{0, 1\}$ be a function which is not identically $0$ with $f(0^n) = 0$,
then $\mathbf{N}(f) = \lceil\log(a(f) + 1)\rceil - 1$.
\end{theorem}

Blais et al.~\cite{BCOST15} showed decomposition for functions with low alternating number~\cite{BCOST15}.
\begin{theorem}[Theorem 1.1 in~\cite{BCOST15}]
Let $f$ be a $k$-alternating function, then $f(x) = h(m_1(x),\dots, m_k(x))$
where $m_i(x)$ is monotone and h is the parity function or its negation. Conversely, any
function of this form is $k$-alternating.
\end{theorem}

The above characterization shows a simple structure for functions with a low alternating number,
which are computable by few negation gates.
To simplify notation, %
we will %
focus on the parity of few monotone functions.

\begin{definition} [$k$-monotone function]
A function $f:\{0,1\}^n\rightarrow\{0,1\}$ is said to be $k$-monotone, if there exist
$k$ monotone functions $g_1, g_2, \dots, g_k$ such that $f = g_1 \oplus g_2 \oplus \cdots \oplus g_k$. 

\end{definition}

\subsection{Orthogonal basis for functions over a slice}
%
%
%
%
Given a subset of the $r$-slice, $A\subseteq \binom{[n]}{r}$, %
denote its density in this slice by $\mu(A)$, \ie, 
%
$\mu(A)=|A|/\binom{n}{r}$. %
%
%
Define its \emph{upper shadow}  %
as
%
\begin{equation} %
%
\partial^+ A:=\left\{x\in
%
\binom{[n]}{r+1} %
:x\succ y \text{ for some }y\in A\right\} %
%
\end{equation} %
%
and its \emph{lower shadow} as  %
%
\begin{equation} %
%
\partial^- A:=\left\{x\in
%
\binom{[n]}{r-1} %
:x\prec y \text{ for some } y\in A\right\}. %
%
\end{equation} %

Filmus~\cite{Filmus} and Srinivasan~\cite{srinivasan2011symmetric} independently 
introduced an orthogonal basis for functions over a
%
slice 
%
$\binom{[n]}{r}$ %
of the Boolean hypercube,   %
which plays a central role in our proofs. 
%
Now we describe the explicit formula provided by Filmus~\cite{Filmus} (see also Section 9.2 in \cite{FilmusM16}). %
%
%
%
%

%
%
%
%
%
%
%
%
%
%
%

%
%
%
%
%
%
%
%
%
%
%
%
%
%
%
%
\begin{definition}
For $d\leq n/2$, define a \emph{$\mathcal{B}_{n,d}$-sequence} to be a sequence $B=b_1,b_2,\cdots,b_d$ where (i) all $b_1,\cdots,b_d$ belong to $[n]$, (ii) $b_1<b_2<\cdots<b_d$, and (iii) $b_i\geq 2i$ for any $1\leq i\leq d$.
We will also use $\mathcal{B}_{n,d}$ to denote the set of all  $\mathcal{B}_{n,d}$-sequences.  

Given $B\in\mathcal{B}_{n,d}$, define $\mathcal{T}(B)$ to be the set consisting of all sequences $A=a_1,a_2,\cdots,a_d$ of $d$ distinct elements of $[n]$ satisfying that (i) $\{a_1,\cdots,a_d\}\cap \{b_1,\cdots,b_d\}=\emptyset$ and (ii) $a_i<b_i$ for all $i$. Furthermore, we define
\begin{equation} 
\chi_B=\sum_{A\in \mathcal{T}(B)} \prod_{i=1}^d(x_{a_i}-x_{b_i}).
\end{equation} %
\end{definition}
%
\begin{theorem}[Theorem 15 in~\cite{Filmus}]\label{thm:def} %
Let $r\leq n/2$ be an integer, the set $\{\chi_B:B\in \mathcal{B}_{n,d}$ for some $d\leq r \}$ is an orthogonal basis for the vector space of functions over the slice 
%
$\binom{[n]}{r}$. %
The Young-Fourier expansion of $f:
%
\binom{[n]}{r} %
\rightarrow \mathbb{R}$ is the unique representation
%
\begin{equation} %
f=\sum_{B\in \mathcal{B}_{n,d},\ d\leq r}\hat{f}(B)\chi_B,
%
\end{equation} %
where $\hat{f}(B)=\frac{\Braket{f,\chi_B}}{\|\chi_{B}\|_2^2}$. Here $\Braket{f,g}:=\mathbb{E}_{x\sim U}[f(x)g(x)]$. In addition for $B\in \mathcal{B}_{n,d}$,
\begin{enumerate}
  \item $\| \chi_{B}\|_2^2=\prod_{i=1}^{d}\frac{(b_i-2(i-1))(b_i-2(i-1)-1)}{2}\cdot 2^d\frac{r^{\underline{d}}(n-r)^{\underline{d}}}{n^{\underline{2d}}}=n^{O(d)}$. In particular, if $r\geq \frac{n}{4}$ and $d=o(n)$, then %
  $\|\chi_{B}\|_2^2=2^{-O(d)}$. %
  Here, $r^{\underline{d}}=\prod_{i=0}^{d-1}(r-i)$.
  %
  \item $\| \chi_{B}\|_{\infty}\leq\sum_{A\in \mathcal{T}(B)}\|\chi_{A,B}\|_{\infty}=n^{O(d)}$. %
\end{enumerate}
\end{theorem}

By Boolean duality, we can extend the above Young-Fourier expansion to where $r>n/2$ naturally. This can be done by replacing the basis $\{\chi_B(x)\}$ by $\{\chi_{\bar{B}}(x):=\chi_B(\bar{x})\}$ where $\bar{x}$ is obtained by flipping all bits of $x$.
\begin{corollary}\label{thm:def2}
    Let $r>n/2$ be an integer, the set $\{\chi_{\bar{B}}(x):B\in \mathcal{B}_{n,d}$ for some $d\leq n-r\}$ is an orthogonal basis for the vector space of functions over the slice 
    %
    $\binom{[n]}{r}$. %
    The Young-Fourier expansion of $f:
    %
    \binom{[n]}{r} %
    \rightarrow \mathbb{R}$ is the unique representation
%
\begin{equation} %
f=\sum_{B\in \mathcal{B}_{n,d},\ d\leq n-r}\hat{f}(B)\chi_{\bar{B}},
%
\end{equation} %
where $\hat{f}(B)=\frac{\Braket{f,\chi_{\bar{B}}}}{\|\chi_{\bar{B}}\|_2^2}$. In addition, we have $\|\chi_{\bar{B}}\|_2^2=\|\chi_{B}\|_2^2$ and $\|\chi_{\bar{B}}\|_\infty=\|\chi_{B}\|_\infty$.
\end{corollary}

Like %
for  %
functions over the Boolean hypercube, we can define the 
%
\emph{total weight} %
on level $d$:
\begin{definition} Let $f:
%
\binom{[n]}{r} %
\rightarrow\mathbb{R}$, for any $r\leq n$, define
%
\begin{equation} %
\W^d(f) = \sum_{B\in \mathcal{B}_{n,d}}\hat{f}(B)^2\| \chi_B\|_2^2,
%
\end{equation} %
and denote $\W^{>d}(f) = \sum_{d' > d} \W^d$ and $\W^{\leq d}(f) = \sum_{d'\leq  d} \W^d$.
\end{definition}
\begin{definition}
Let $f:
%
\binom{[n]}{r} %
\rightarrow \{\pm1\}$. For $i,j\in [n]$, define the 
%
\emph{influence} %
of $f$ on the pair $(i,j)$ as
%
\begin{equation} %
\I_{ij}[f]=2\Pr[f(x^{(i,j)})\neq f(x)].
%
\end{equation} %
Here $x^{(i,j)}$ is obtained by switching $x_i$ and $x_j$.
The %
\emph{total influence} %
of $f$ is
%
\begin{equation} %
\I[f]=\frac{1}{n}\sum_{1\leq i<j\leq n}\I_{ij}[f].
%
\end{equation} %
\end{definition}
%
\begin{lemma}[Proposition 3.1 in~\cite{kkl09}]\label{key} %
    Let $f:
    %
    \binom{[n]}{r} %
    \rightarrow \{\pm1\}$ and $A=\{x\in 
    %
    \binom{[n]}{r}: %
    f(x)=-1\}$, then
    %
\begin{equation} %
\min\{\mu(\partial^+ A),\mu(\partial^- A)\}\geq \mu(A)+\frac{n}{4r(n-r)}\cdot \I[f]\geq\mu(A)+\frac{1}{n}\I[f].
    %
\end{equation} %
\end{lemma}

%
\begin{theorem}[Lemma 24 in~\cite{Filmus}]\label{slice_inf} %
Let $f:
%
\binom{[n]}{r} %
\rightarrow \{\pm1\}$. Then
%
\begin{equation} %
\I[f]=\sum_{d \leq \min(r,n-r)}\frac{d(n+1-d)}{n} \cdot \W^d
 = \sum_{B\in \mathcal{B}_{n,d}, d \leq \min(r,n-r)}\frac{d(n+1-d)}{n} \cdot \hat{f}(B)^2\| \chi_B\|_2^2.
%
\end{equation} %

In addition, according to Parseval's identity,
%
\begin{equation} %
\sum_d \W^d = \sum_{B\in \mathcal{B}_{n,d},d\leq \min(r,n-r)}\hat{f}(B)^2\| \chi_B\|_2^2=\| f\|_2^2=1.
%
\end{equation} %
\end{theorem}

\subsection{Basic inequalities}
Finally, we will make use of the Hoeffding bound.
%
%
\begin{theorem}[Hoeffding Bound, Theorem 2 in \cite{Hoeffdingbound}] %
	Let $X=\sum_{i=1}^{n}X_i$, where $X_i\in[a_i,b_i]$ are independent random variables. Then for any $\theta>0$,
	%
\begin{equation} %
\Pr(|X-\mathbb{E}(X)|\geq \theta)\leq 2\exp\left(-\frac{2\theta^2}{\Sigma_{i}(b_i-a_i)^2}\right).
%
\end{equation} %
\end{theorem}

\begin{corollary}\label{Coro:Heff}
  Let $X$ be a random variable with distribution $\mathcal{D}$ whose range is $[l,u]$. Let $X_1,\ldots,X_m$ be its independent samples. Then w.p.
%
$\ge 1-\delta$,  %
for any $\epsilon>0$,
    %
\begin{equation} %
\abs{\frac{1}{m}\sum_{i=1}^{m}X_i-\mathbb{E}(X)}\leq \epsilon
%
\end{equation} %
	as long as $m\geq (u-l)^2\log(2/\delta)/(2\epsilon^2)$.
\end{corollary}

\warn{
The following fact will also be used.
\begin{proposition}\label{prop:binomial}
For
%
%
$t=o(n)$ we have  %
%
$\binom{n}{n/2-t/2}/2^n %
=\frac{1}{\sqrt{n}}\cdot 2^{-\Theta(t^2/n)}$.
%
\end{proposition}
\begin{proof}
By Stirling's approximation,
%
\begin{equation} %
\log %
\binom{n}{n/2-t/2} %
=\log(1+o(1))+\log\sqrt{\frac{n}{2\pi(n/2-t/2)(n/2+t/2)}}+n\cdot H\left(\frac{n/2-t/2}{n}\right),
%
\end{equation} %
where $H(p)=-p\log p-(1-p)\log(1-p)$ is the binary entropy function.
As $t/2n=o(1)$, by the Taylor expansion of the entropy function around $1/2$, we have 
%
\begin{equation} %
H\left(\frac{1}{2}-\frac{t}{2n}\right)=1-\frac{1+o(1)}{2\ln 2}\left(\frac{t}{n}\right)^2.
%
\end{equation} %
The conclusion follows immediately.
\end{proof}
}


\section{Concentration property of $k$-monotone functions}
\warn{In the rest of this paper, for a function $f:\{0,1\}^n\rightarrow\{0,1\}$, we convert the range to $\{\pm 1\}$. 
The mapping from $\{0,1\}$ to $\{-1,1\}$ is given by $1-2b$, sending $0$ to $1$ and $1$ to $-1$. So a function $f:\{0,1\}^n\rightarrow\{\pm 1\}$ is said to be $k$-monotone if $(1-f)/2$ is $k$-monotone.}

In this section, we show 
that %
some ``middle'' slice
of a $k$-monotone function has Fourier concentration. 
For functions $f:\{0,1\}^n\rightarrow\{\pm1\}$, let $f|_r$ be the subfunction of $f$ restricted to %
$\binom{[n]}{r}$ %
and $\mu(f|_r):=\mu(f|_r^{-1}(-1))$.

\begin{definition}[$(t,d,\epsilon)$-concentration]
We say $f:\{0,1\}^n\rightarrow\{\pm1\}$ is $(t,d,\epsilon)$-concentrated if the following holds: for some $r$ such that $n/2-t/2\leq r\leq n/2+t/2$,
%
\begin{equation} %
\W^{>d}(f|_r)=\sum_{B\in \mathcal{B}_{n,d'}: d'> d }\widehat{f|_r}(B)^2\| \chi_B\|_2^2 < \epsilon.
%
\end{equation} %
\warn{Intuitively, $f$ has low-degree Fourier concentration on at least one of the middle slices.}
\end{definition}

\begin{lemma} \label{lm:main}
Let $f:\{0,1\}^n\rightarrow\{\pm1\}$ be a $k$-monotone function.
For any $1<t \leq n$ and any $d$, $\epsilon$ such that $d\epsilon \geq 2kn/t$, 
we have that %
$f$ is $(t, d, \epsilon)$-concentrated.
\end{lemma}

\expref{Lemma}{lm:main} follows from an upper bound on the sum of total influences on slices.
\begin{proposition}\label{prop:main}
Let $f:\{0,1\}^n\rightarrow\{\pm1\}$ be a $k$-monotone function. Then $\sum_{r=0}^{n-1} \I[f|_r]\leq kn$.
\end{proposition}


We first prove \expref{Lemma}{lm:main} using \expref{Proposition}{prop:main}.

\begin{proof}[Proof of Lemma \ref{lm:main}]
By contradiction, we assume that $\W^{>d}(f|_r)>\epsilon\geq \frac{2kn}{dt}$ for any $n/2-t/2\leq r\leq n/2+t/2$.
  According to \expref{Proposition}{prop:main},  we have $\sum_{r=\lceil n/2-t/2\rceil}^{\floor{n/2+t/2}} \I[f|_r]\leq kn$. By averaging, let $n/2 - t/2\leq r\leq n/2 + t/2$ be such that
  $\I[f|_r] \le kn/t$.
  By \expref{Theorem}{slice_inf}, we can deduce that,
  %
  \begin{align} %
    \frac{kn}{t} \ge \I[f|_r] & =\sum_{d' \le \min(r,n-r)}\frac{d'(n+1-d')}{n} \cdot \W^{d'}(f|_r)\\
     & \ge \sum_{d < d' \le \min(r,n-r)}\frac{d'(n+1-d')}{n}\cdot \W^{d'}(f|_r) \\
     & \ge \frac{d(n+1-d)}{n} \cdot \sum_{d < d' \le \min(r, n-r)}\W^{d'}(f|_r) \\
     & > \frac{d}{2} \cdot \epsilon\geq \frac{d}{2} \cdot \frac{2kn}{dt}\geq \frac{kn}{t} \; ,
  %
  \end{align} %
a contradiction.
\end{proof}

Now we prove \expref{Proposition}{prop:main}.

\begin{proof}[Proof of Proposition \ref{prop:main}]
    Suppose $f$ is the parity of $h_1,\cdots,h_k$ where each $h_i$ is monotone.
    For any $r$, when we switch $x_i$ and $x_j$, $f|_r(x)$ changes only if at least one $h_i|_r(x)$ changes for $i=1,2,\dots, k$. Thus, combining with the union bound, we have
    \begin{align}\label{eq:sec3_1}
    \I[f|_r]\leq \sum_{i=1}^k\I[h_i|_r].
    \end{align}
    
    \warn{Since $h_i$ is monotone, the upper shadow of $h_i|^{-1}_r(-1)$ is a subset of $h_i|^{-1}_{r+1}(-1)$. Then according to \expref{Lemma}{key}, we have
    %
\begin{equation} %
    \mu(h_i|_{r+1})\geq \mu(h_i|_r)+\frac{1}{n}I(h_i|_r),
    %
\end{equation} %
    which implies 
    \begin{align}\label{eq:sec3_2}
    \frac{1}{n}\sum_{r=0}^{n-1} \I[h_i|_r]\leq \mu(h_i|_{n})-\mu(h_i|_{0})\leq 1.
    \end{align}
    
    \expeqref{Equation}{eq:sec3_1} and \expeqref{Equation}{eq:sec3_2} imply the desired conclusion.}
\end{proof}
 
\section{Distinguishers for $k$-monotone functions}

In this section, we prove the following theorem.

\distinguisher*

We prove this theorem by giving a distinguisher for $(t,d,1/2)$-concentrated functions.

\begin{proposition}
For $t\leq n/4$ and $d=o(n/\log n)$, any family of $(t,d,1/2)$-concentrated functions can be distinguished from uniform random functions in
$2^{O(d\log n + t^2/n)}$ time.
\end{proposition}

By \expref{Lemma}{lm:main}, every $k$-monotone function is \warn{ $\left((kn^2\log{n})^{1/3},4\big(\frac{k^2n}{\log n}\big)^{1/3},1/2\right)$}-concentrated, then \expref{Theorem}{thm:distinguisher} follows. Now we prove the proposition.

\begin{proof}
The distinguisher is given in \expref{Algorithm}{alg:distinguisher}. %
We will %
show that
\begin{itemize}
    \item ({\bf Soundness}) It accepts a uniform random function w.p. $o(1)$;
    \item ({\bf Completeness}) It accepts any $(t,d,1/2)$-concentrated function w.p. $1-o(1)$;
    \item ({\bf Complexity}) Its sample/time complexity is $2^{O(d\log n + t^2/n)}$.
\end{itemize}

\begin{algorithm}[t]
\caption{A Distinguisher for $(t,d,1/2)$-Concentrated Functions}\label{alg:distinguisher}

    Let $C$ be a large enough constant\;
    \For{$r\leftarrow \lceil\frac{n}{2}-\frac{t}{2}\rceil$ \KwTo $\lfloor\frac{n}{2}+\frac{t}{2}\rfloor$}
 {
     $S\leftarrow 0$\;
     \For{$B\in \mathcal{B}_{n,d'}$ with $d'\leq d$}
     {
          \If {$r\leq n/2$}
           {Estimate $\Braket{f|_r, \chi_B}$ with accuracy $n^{-C\cdot d}$\;
           $S\leftarrow S+\widehat{f|_r}(B)^2\| \chi_B\|_2^2$\;
          \tcp{$\widehat{f|_r}(B)^2\| \chi_B\|_2^2 = \Braket{f|_r, \chi_B}^2/\|\chi_B\|_2^2$}}
          \Else {Estimate $\Braket{f|_r, \chi_{\bar{B}}}$ with accuracy $n^{-C\cdot d}$\;
           $S\leftarrow S+\widehat{f|_r}(\bar{B})^2\| \chi_{\bar{B}}\|_2^2$\;
          \tcp{$\widehat{f|_r}(B)^2\| \chi_{\bar{B}}\|_2^2= \Braket{f|_r, \chi_{\bar{B}}}^2/\|\chi_{\bar{B}}\|_2^2$}}
     }
     \If{$S\geq 3/8$}
        {Return True\;}
 }
 Return False\;

\end{algorithm}

\vspace{2ex}
\noindent\textbf{Soundness.}\quad Let $f$ be a uniform random function. We claim that for each $\frac{n}{2}-\frac{t}{2}\leq r\leq \frac{n}{2}+\frac{t}{2}$, the variable $S$ in Line 11 is at most $1/4$ w.p. $1-o(\frac{1}{n})$, which concludes the soundness by the union bound. 

Fix such an $r$. W.l.o.g., we assume that $r\leq n/2$. For any $B\in\mathcal{B}_{n,d'}$ where $d'\leq d$, it is easily seen that $\E_f\big[\Braket{f|_r, \chi_B}\big]=0$, then by the Hoeffding bound,
%
\begin{equation} %
    \Pr_f\left[\left|\Braket{f|_r, \chi_B}\right|\ge \theta\right]\leq 2 \exp\left(-2\theta^2 
    %
    \binom{n}{r} %
    \big/ \|\chi_B\|^2_2\right)\leq 2\exp\left(-\theta^2 \left(\frac{4}{3}\right)^{n/4-o(n)}\right),
%
\end{equation} %
where the last inequality is due to that %
$\binom{n}{r} %
\geq (\frac{n}{r})^r \geq (\frac{4}{3})^{n/4}$ and $\|\chi_B\|^2_2=2^{O(d\log n)}=(\frac{4}{3})^{o(n)}$.
In particular, by letting $\theta=\left(\frac{3}{4}\right)^{n/10}$ and using the union bound, we have that with probability at least $1-n^{O(d)}\cdot \exp\left(-\theta^2 \left(\frac{4}{3}\right)^{n/4-o(n)}\right)=1-o(\frac{1}{n})$, $\big|\Braket{f|_r, \chi_B}\big|< \theta$ for every $B\in\mathcal{B}_{n,d'}$ where $d'\leq d$. Thus, w.p. $1-o(\frac{1}{n})$,
%
\begin{align} %
    \W^{\leq d}[f|_r]&=\sum_{B\in\mathcal{B}_{n,d'},d'\leq d} \widehat{f|_r}(B)^2\|\chi_B\|_2^2=\sum_{B\in\mathcal{B}_{n,d'},d'\leq d} \frac{\Braket{f|_r, \chi_B}^2}{\|\chi_B\|_2^2}
 \leq \sum_{B\in\mathcal{B}_{n,d'},d'\leq d} \frac{\theta^2}{\|\chi_B\|_2^2}\leq \frac{1}{8},
%
\end{align} %
where the last inequality holds for sufficiently large $n$. Finally, $S$ is an estimate of $\W^{\leq d}[f|_r]$ with additive error $n^{-\Omega(d)}$.

\vspace{2ex}
\noindent\textbf{Completeness.}\quad Let $f$ be a $(t,d,1/2)$-concentrated function. By definition, there is some $r$ such that $n/2-t/2\leq r\leq n/2+t/2$ and $W^{\leq d}[f|_r]>1/2$. As $S$ is an estimate of $\W^{\leq d}[f|_r]$ with additive error $n^{-\Omega(d)}$, we conclude that \expref{Algorithm}{alg:distinguisher} accepts $f$ with high probability.

\vspace{2ex}
\noindent\textbf{Complexity.}\quad The loop in Line~2 is repeated at most $t$ times. In Line~4, the number of strings 
%
$B\in \mathcal{B}_{n,d'}$ with $d'\leq d$ %
we enumerated is at most $n^{O(d)}$.
Furthermore, for each $n/2-t/2\leq r\leq n/2+t/2$ and each 
%
$B\in\mathcal{B}_{n,d'}$ with $d'\leq d$, 
according to the Hoeffding bound, $n^{O(d)}$ uniform random samples on the slice 
%
$\binom{[n]}{r}$ %
are sufficient to estimate $\Braket{f|_r, \chi_B}$ with accuracy $n^{-C\cdot d}$. In addition, a random uniform sample is from the slice 
%
$\binom{[n]}{r}$ %
with probability 
%
$\binom{n}{r}/2^n$, %
which is $\frac{1}{\sqrt{n}}\cdot2^{-O(t^2/n)}$ according to \expref{Proposition}{prop:binomial}. Thus, the total number of random samples used is at most $t \cdot n^{O(d)}\cdot n^{O(d)}\cdot 2^{O(t^2/n)}=2^{O(d\log n+t^2/n)}$.

Besides, the function 
%
$\chi_B=\sum_{A\in \mathcal{T}(B)}\chi_{A,B}$ %
can be computed by enumerating all $n^{O(d)}$ strings $A$ in 
%
$\mathcal{T}(B)$. %
Thus, the time complexity is also $2^{O(d\log n + t^2/n)}$.
\end{proof}


%
\begin{algorithm}[t]
\caption{A weak learner for $(t,d, 1/2)$-concentrated functions}
\label{alg:learner}
  Let $C$ be a large enough constant\;
 \For{$r\leftarrow \frac{n}{2}-\frac{t}{2}$ \KwTo $\frac{n}{2}+\frac{t}{2}$}
 {
     $S\leftarrow 0$\;
     $p\leftarrow 1$\;
     \For{$B\in \mathcal{B}_{n,d'}$ with $d'\leq d$}
     {
        \If{$r\leq n/2$}
        {
            Estimate $\Braket{f|_r, \chi_B}$ with accuracy $n^{-C\cdot d}$\;
            $S\leftarrow S+\widehat{f|_r}(B)^2\| \chi_B\|_2^2$\;
        }\Else
        {
            Estimate $\Braket{f|_r, \chi_{\bar{B}}}$ with accuracy $n^{-C\cdot d}$\;
            $S\leftarrow S+\widehat{f|_r}(B)^2\| \chi_{\bar{B}}\|_2^2$\;
        }
          
     }
     \If{$S\geq 3/8$}
       {
       \If{$r\leq n/2$}{
            $g(x)\leftarrow \sum_{B\in\mathcal{B}_{n,d'}:d'\leq d}\widehat{f|_r}(B)\chi_B(x)$\; 
       }\Else{
            $g(x)\leftarrow \sum_{B\in\mathcal{B}_{n,d'}:d'\leq d}\widehat{f|_r}(B)\chi_{\bar{B}}(x)$\;
       }

        \While{$p>\sqrt{3}/4$}
        {
            Pick $\theta \in [-1,1]$ uniformly at random\;
            Estimate $p\leftarrow\Pr[f|_r\neq\operatorname{sign}(g-\theta)]$\;
        }
        Estimate $\mu_{\neq r}\leftarrow \mathbb{E}_{x}[f(x)\mid |x|\neq r]$\;

        Return $h(x)=\begin{cases}
                         \operatorname{sign}(g(x)-\theta) & \mbox{if } |x|=r;\\
                        \operatorname{sign}(\mu_{\neq r}) & \mbox{if } |x|\neq r.
                     \end{cases}$
        }
  
 }
 Return $h(x)\equiv 0$.
\end{algorithm}

\section{Weak learners for $k$-monotone functions}
In this section, we prove the following theorem.

\learner*

We prove \expref{Theorem}{thm:learner} by giving a weak learner for $(t,d,1/2)$-concentrated functions. By \expref{Lemma}{lm:main}, $k$-monotone functions are $\left(\sqrt{n\log n},4k\sqrt{\frac{n}{\log n}},1/2\right)$-concentrated, then \expref{Theorem}{thm:learner} follows.

\begin{proposition}\label{prop:learner}
    For $t=O(\sqrt{n\log n})$ and $d=o(n/\log n)$, \expref{Algorithm}{alg:learner} weakly learns $(t,d,1/2)$-concentrated functions in $2^{O(d\log n + t^2/n)}$ time.
\end{proposition}

To learn $(t,d,1/2)$-concentrated functions $f$, \expref{Algorithm}{alg:learner} tries to find out the slice %
$\binom{[n]}{r}$ %
on which $f|_r$ is concentrated, and then figures out a function $g:
%
\binom{[n]}{r} %
\rightarrow \mathbb{R}$ which is very close to $f|_r$. To convert the approximated function $g$ to a Boolean-valued function, we can utilize \expref{Claim}{lem:learn} similar to Exercise 3.34 in \cite{book}.
For the rest of the slices, the learner just outputs the most frequent value. \warn{Since $t=O(\sqrt{n\log n})$, each slice in $[n/2-t/2, n/2+t/2]$ is at least a $\frac{1}{\sqrt{n}}\cdot 2^{-O(t^2/n)}=1/\poly(n)$ fraction according to \expref{Proposition}{prop:binomial}. Hence we get a $(1/2-1/\poly(n))$-close function $h$.}

\begin{proof}[Proof of Proposition~\ref{prop:learner}]
We first show that \expref{Algorithm}{alg:learner} weakly learns $(t,d,1/2)$-concentrated functions. Let $f$ be a $(t,d,1/2)$-concentrated function. For each $n/2-t/2\leq r\leq n/2+t/2$, the variable $S$ in Line 12 is an estimate of $\W^{\leq d}[f|_r]$ with additive error $n^{-\Omega(d)}$. Then for some $r$, the condition $S\geq 3/8$ in Line 12 holds, and \expref{Algorithm}{alg:learner} executes Lines 13-21. 
For the function $g$ obtained in Line 14 or Line 16, and sufficiently large $n$, 
\begin{align*}
  ||f|_{r}-g||_1\leq ||f|_{r}-g||_2 = \sqrt{\sum_{B\in \mathcal{B}_{n,d'}, d'\leq d}{(\widehat{f|_{r}}(B)-\hat{g}(B))^2}\|\chi_B\|_2^2+W^{>d}[f|_{r}]}\leq \sqrt{o(1)+5/8}\leq \sqrt{3}/2.
\end{align*}
To convert $g$ to a Boolean-valued function, we utilize the following claim. 
\begin{claim}\label{lem:learn}
 Suppose $f:
 %
 \binom{[n]}{r} %
 \rightarrow\{-1,1\}$ and $g:
 %
 \binom{[n]}{r} %
 \rightarrow \mathbb{R}$. Pick $\theta \in [-1,1]$ uniformly at random and define $g'=\operatorname{sign}\big(g(x)-\theta\big)$, we have $\mathbb{E}_\theta\left[\Pr_x\big(f(x)\neq g'(x)\big)\right]\leq \|f-g\|_1/2$.
\end{claim}
\begin{proof} By rewriting the last formula and swapping the expectation operators, we have
\begin{equation*}
     \begin{aligned}
        \E_\theta \left[\Pr_x\big(f(x)\neq g'(x)\big)\right] &= \E_\theta \E_x \left[1_{f(x) \neq g'(x)}\right] = \E_x\E_\theta \left[1_{f(x) \neq g'(x)}\right]
        = \E_x\left[\Pr_\theta \big(f(x) \neq \operatorname{sign}(g(x)-\theta)\big)\right]\\
          &\leq\E_x\left[\frac{|f(x)-g(x)|}{2}\right]=\frac{\|f-g\|_1}{2}. & \qedhere
%
    \end{aligned}
\end{equation*}
\end{proof}
Thus, for a random $\theta\in[-1,1]$, $\Pr\big[f|_{r}(x)\neq \operatorname{sign}(g(x)-\theta)\big] \leq \sqrt{3}/4$ holds with a constant probability. That is, with high probability, the loop of Lines 17-19 is repeated a constant number of times, and we will get a $\theta^\star$ such that $\Pr\big[f|_{r}(x)\neq \operatorname{sign}(g(x)-\theta^\star)\big] \leq \sqrt{3}/4$. Finally, we have
%
%
\begin{align} %
& \Pr[h(x)\neq f(x)] \\
= &\Pr[|x|\neq r]\Pr[f(x)\neq \operatorname{sign}(\mu_{\neq r})\mid |x|\neq r]+\Pr[|x|=r]\Pr\big[f|_{r}(x)\neq \operatorname{sign}(g(x)-\theta^\star)\big]\\
\leq &\left(1-
%
\frac{\binom{n}{r}}{2^n} %
\right)\cdot \frac{1}{2}+ 
%
\frac{\binom{n}{r}}{2^n}  %
\cdot \frac{\sqrt{3}}{4} =\frac{1}{2} + 
%
\frac{\binom{n}{r}}{2^n} %
\cdot\left(\frac{\sqrt{3}}{4}-\frac{1}{2}\right)\\=&\frac{1}{2}-\left(\frac{1}{2}-\frac{\sqrt{3}}{4}\right)\cdot\frac{1}{\sqrt{n}}\cdot 2^{-O(t^2/n)}=\frac{1}{2} - \frac{1}{\mbox{poly}(n)},
%
\end{align} %
\warn{where the
%
penultimate  %
equality is according to \expref{Proposition}{prop:binomial} and the last equality is due to %
the assumption that  %
$t=O(\sqrt{n\log n})$.}

What remains is to show that \expref{Algorithm}{alg:learner} terminates in $2^{O(d\log n + t^2/n)}$ time. First, as shown in the analysis of \expref{Algorithm}{alg:distinguisher}, for each $n/2-t/2\leq r\leq n/2+t/2$, it costs $2^{O(d\log n + t^2/n)}$ time to execute Lines 5-11. For some $r$, \expref{Algorithm}{alg:learner} would execute Lines 13-21. As shown above, the loop of Lines 17-19 is  repeated a constant of times. So, it costs $2^{O(d\log n + t^2/n)}$ time to execute Lines 13-21. Therefore the total time complexity is $t\cdot 2^{O(d\log n + t^2/n)}=2^{O(d\log n + t^2/n)}$.
\end{proof}

\section{Discussion and open problems}\label{sec:discussion}

\noindent\textbf{Fourier analysis on slices.}\quad It is surprising to us that a simple variant of the ``Low-Degree Algorithm'' on slices can outperform the classic ``Low-Degree Algorithm'' in terms of attacking negation-limited weak PRFs. To the best of our knowledge, unlike Fourier analysis on the Boolean cube, Fourier analysis on slices has not been explored in cryptography. It is
%
a very  %
interesting direction to use this technique to attack more cryptographic constructions, particularly ones which are secure against attacks based on standard Fourier analysis. 

\vspace{2ex}
%
\noindent\textbf{The hardness of $1$-negation weak PRFs.} One of the most intriguing open problems is how hard can $1$-negation weak PRFs be? Our bound suggests that, unlike testing $1$ negation (using a 1-sided non-adaptive tester)~\cite{GKW19} and learning $1$ negation to high accuracy~\cite{BCOST15}, distinguishing $1$ negation is significantly more efficient than $2^{O(\sqrt{n})}$. Can we have polynomial time distinguishers? We believe that new structural results of $2$-monotone functions are required for polynomial time distinguishers.  

\vspace{2ex}
\noindent\textbf{Fourier spectrum of $k$-monotone functions on low levels.} It is a basic fact \cite{book} that every monotone function $f:\{0,1\}^n\rightarrow\{-1,1\}$ has a large Fourier coefficient on the first two levels. Does a similar statement hold for $k$-monotone functions? We are particularly interested in and focus on the case when $k=o(n)$. 
On one hand, there exist $k$-monotone functions $f$ (e.g., $f=x_1\oplus\cdots\oplus x_k$) such that $\max_{|S|\leq k-1}|\hat{f}(S)|=0$. Second, the minimum $\max_{|S|\leq k}|\hat{f}(S)|$ among all $k$-monotone functions that we are aware of so far is about $\left(\frac{\ln (n/k)}{n/k}\right)^k=\left(\frac{k}{n}\right)^{O(k)}$, which is achieved by the parity of $k$ $\text{Tribes}_{n/k}$ functions on disjoint variables. So we %
are curious about the following conjectures. %

%
%
%
%
%
\begin{conjecture}[Rocco Servedio, 2014]  %
\label{conj:fourier_spectrum} 
For $k=o(n)$ and any $k$-monotone function $f:\{0,1\}^n\rightarrow \{-1,1\}$, there exists a set $S\subseteq[n]$ of size at most $k$ such that $|\hat{f}(S)|=\left(\frac{k}{n}\right)^{O(k)}$.
\end{conjecture}

If \expref{Conjecture}{conj:fourier_spectrum} is true, then applying ``Low-Degree Algorithm'' on the Boolean cube~\cite{LMN93}, we can distinguish $k$-monotone functions from uniformly random functions in $(n/k)^{O(k)}$ time. When $k$ is constant, it also leads to a polynomial time weak learner. We conjecture $(n/k)^{\Theta(k)}$ is the correct bound of distinguishers, specifically, there exists a $\left((n/k)^{\Theta(k)},1/3\right)$-secure $k$-monotone weak pseudorandom family.
%


\iffalse
\begin{conjecture}\label{conj:weak_prfs}
For $k=o(n)$, there exists a family of $k$-monotone function $f:\{0,1\}^n\rightarrow \{-1,1\}$ such that is $\left((n/k)^{\Theta(k)},1/3\right)$-secure weak pseudorandom family.
\end{conjecture}
\fi


In fact, our first attempt to distinguish $k$-monotone functions is to prove \expref{Conjecture}{conj:fourier_spectrum}. So far, even the following much weaker conjecture remains open.
%
%
%
%
\begin{conjecture}[Rocco Servedio, 2014]  %
\label{conj:fourier_spectrum2} 
	Let $f:\{0,1\}^n\rightarrow \{-1,1\}$ be a $2$-monotone function. There exists a set $S\subseteq[n]$ of size $o(\sqrt{n})$ such that $|\hat{f}(S)|>0$.
\end{conjecture}

We note that it is easy to derive that there exists a set $S\subseteq[n]$ of size $O(\sqrt{n})$ such that $|\hat{f}(S)|>0$ by the $O(\sqrt{n})$ total influence upper bound of $2$-monotone function shown by Blais et al~\cite{BCOST15}.  We are not aware of other relevant results or implications.  %


\vspace{4ex}
\noindent\textbf{Acknowledgments.} We sincerely thank the
ToC editors and the   %
anonymous ToC reviewers 
%
for their detailed and constructive comments. We thank Shengyu Zhang for fruitful discussions at the early stage of this work. Siyao Guo would like to thank Igor Carboni Oliveira for telling her \expref{Conjecture}{conj:fourier_spectrum2} and an early version of \expref{Conjecture}{conj:fourier_spectrum}. 


%
%


%% Theory of Computing
%% File: v020a002.bbl
%% Publication date: July 29, 2024
%% Authors: Zhihuai Chen, Siyao Guo, Qian Li, Chengyu Lin, and Xiaoming Sun
\newcommand{\prelim}{Preliminary version in\ }\newcommand{\prelims}{Preliminary
  versions in\ }
% Define empty bibhead if not already defined
\providecommand{\bibhead}[1]{}
% Define tocrefpdfbookmark if not already defined
\expandafter\ifx\csname pdfbookmark\endcsname\relax%
  \providecommand{\tocrefpdfbookmark}{}
\else\providecommand{\tocrefpdfbookmark}{%
   \hypertarget{tocreferences}{}%
   \pdfbookmark[1]{References}{tocreferences}}%
\fi

\tocrefpdfbookmark
\begin{thebibliography}{10}

\bibitem{ABGKM14}\bibhead{ABGKM14}
{\sc Adi Akavia, Andrej Bogdanov, Siyao Guo, Akshay Kamath, and Alon Rosen}:
  Candidate weak pseudorandom functions in {$AC^0\circ MOD_2$}.
\newblock In {\em Proc. 5th Innovations in Theoret. Comp. Sci. Conf.
  (ITCS'14)}, pp. 251--260. Schloss Dagstuhl--Leibniz-Zentrum fuer Informatik,
  2014.
\newblock [\epfmtdoi{10.1145/2554797.2554821}]

\bibitem{alon1987monotone}\bibhead{alon1987monotone}
{\sc Noga Alon and Ravi~B Boppana}: The monotone circuit complexity of
  {B}oolean functions.
\newblock {\em Combinatorica}, 7(1):1--22, 1987.
\newblock [\epfmtdoi{10.1007/BF02579196}]

\bibitem{andreev1985method}\bibhead{andreev1985method}
{\sc Alexander~E Andreev}: A method for obtaining efficient lower bounds for
  monotone complexity.
\newblock {\em Algebra and Logic}, 26(1):1--18, 1987.
\newblock \prelim Doklady 282, 1985, pp.1033-1037. Original article in
  \href{https://www.mathnet.ru/php/archive.phtml?wshow=paper&jrnid=al&paperid=1967}{\emph{Algebra
  i Logika} 26(1) 3--26, 1987 (Russian)}.
\newblock [\epfmtdoi{10.1007/BF01978380}]

\bibitem{berg1999symmetric}\bibhead{berg1999symmetric}
{\sc Christer Berg and Staffan Ulfberg}: Symmetric approximation arguments for
  monotone lower bounds without sunflowers.
\newblock {\em Comput. Complexity}, 8(1):1--20, 1999.
\newblock [\epfmtdoi{10.1007/s000370050017}]

\bibitem{BCOST15}\bibhead{BCOST15}
{\sc Eric Blais, Cl{\'{e}}ment~L. Canonne, Igor~Carboni Oliveira, Rocco~A.
  Servedio, and Li{-}Yang Tan}: Learning circuits with few negations.
\newblock In {\em Proc. 18th Internat. Workshop on Approximation Algorithms for
  Combinat. Opt. Probl. (APPROX'15)}, pp. 512--527. Schloss
  Dagstuhl--Leibniz-Zentrum fuer Informatik, 2015.
\newblock [\epfmtdoi{10.4230/LIPIcs.APPROX-RANDOM.2015.512}]

\bibitem{BABL98}\bibhead{BABL98}
{\sc Avrim Blum, Carl Burch, and John Langford}: On learning monotone {B}oolean
  functions.
\newblock In {\em Proc. 39th FOCS}, pp. 408--415. IEEE Comp. Soc., 1998.
\newblock [\epfmtdoi{10.1109/SFCS.1998.743491}]

\bibitem{BR17}\bibhead{BR17}
{\sc Andrej Bogdanov and Alon Rosen}: Pseudorandom functions: Three decades
  later.
\newblock In {\sc Yehuda Lindell}, editor, {\em Tutorials on the Foundations of
  Cryptography}, pp. 79--158. Springer, 2017.
\newblock [\epfmtdoi{10.1007/978-3-319-57048-8\_3}]

\bibitem{CGGKW17}\bibhead{CGGKW17}
{\sc Cl{\'{e}}ment~L. Canonne, Elena Grigorescu, Siyao Guo, Akash Kumar, and
  Karl Wimmer}: Testing $k$-monotonicity: {T}he rise and fall of {B}oolean
  functions.
\newblock {\em Theory of Computing}, 15(1):1--55, 2019.
\newblock \prelim \href{https://doi.org/10.4230/LIPIcs.ITCS.2017.29}{ITCS'17}.
\newblock [\epfmtdoi{10.4086/toc.2019.v015a001}]

\bibitem{DFTWW14}\bibhead{DFTWW14}
{\sc Dana Dachman-Soled, Vitaly Feldman, Li-Yang Tan, Andrew Wan, and Karl
  Wimmer}: Approximate resilience, monotonicity, and the complexity of agnostic
  learning.
\newblock In {\em Proc. 26th Ann. ACM--SIAM Symp. on Discrete Algorithms
  (SODA'15)}, pp. 498--511. SIAM, 2015.
\newblock [\epfmtdoi{10.1137/1.9781611973730.34}]

\bibitem{Filmus}\bibhead{Filmus}
{\sc Yuval Filmus}: An orthogonal basis for functions over a slice of the
  {B}oolean hypercube.
\newblock {\em Electr. J. Combinat.}, 23(1):P1.23, 2016.
\newblock [\epfmtdoi{10.37236/4567}]

\bibitem{FilmusM16}\bibhead{FilmusM16}
{\sc Yuval Filmus and Elchanan Mossel}: Harmonicity and invariance on slices of
  the {B}oolean cube.
\newblock In {\em Proc. 31st Comput. Complexity Conf. (CCC'16)}, pp. 16:1--13.
  Schloss Dagstuhl--Leibniz-Zentrum fuer Informatik, 2016.
\newblock [\epfmtdoi{10.4230/LIPIcs.CCC.2016.16}]

\bibitem{GGM}\bibhead{GGM}
{\sc Oded Goldreich, Shafi Goldwasser, and Silvio Micali}: How to construct
  random functions.
\newblock {\em J. ACM}, 33(4):792--807, 1986.
\newblock [\epfmtdoi{10.1145/6490.6503}]

\bibitem{GKW19}\bibhead{GKW19}
{\sc Elena Grigorescu, Akash Kumar, and Karl Wimmer}: Flipping out with many
  flips: Hardness of testing {$k$}-monotonicity.
\newblock {\em SIAM J. Comput.}, 33(4):2111--2125, 2019.
\newblock [\epfmtdoi{10.1137/18M1217978}]

\bibitem{GK15}\bibhead{GK15}
{\sc Siyao Guo and Ilan Komargodski}: Negation-limited formulas.
\newblock {\em Theoret. Comput. Sci.}, 660:75--85, 2017.
\newblock [\epfmtdoi{10.1016/j.tcs.2016.11.027}]

\bibitem{GMOR15}\bibhead{GMOR15}
{\sc Siyao Guo, Tal Malkin, Igor~C Oliveira, and Alon Rosen}: The power of
  negations in cryptography.
\newblock In {\em Proc. Theory of Cryptography Conf. (TCC'15)}, pp. 36--65.
  Springer, 2015.
\newblock [\epfmtdoi{10.1007/978-3-662-46494-6\_3}]

\bibitem{harnik2000higher}\bibhead{harnik2000higher}
{\sc Danny Harnik and Ran Raz}: Higher lower bounds on monotone size.
\newblock In {\em Proc. 32nd STOC}, pp. 378--387. ACM Press, 2000.
\newblock [\epfmtdoi{10.1145/335305.335349}]

\bibitem{Hoeffdingbound}\bibhead{Hoeffdingbound}
{\sc Wassily Hoeffding}: Probability inequalities for sums of bounded random
  variables.
\newblock {\em J. Amer. Statistical Assoc.}, 58(301):13--30, 1963.
\newblock Also available in
  \href{https://doi.org/10.1007/978-1-4612-0865-5\_26}{The Collected Works of
  Wassily Hoeffding, 1994, pp. 409--426, Springer} and at
  \href{https://doi.org/10.2307/2282952}{JSTOR}. See also Theorem 1 in
  \href{https://web.eecs.umich.edu/~cscott/past\_courses/eecs598w14/notes/03\_hoeffding.pdf}{Clayton
  Scott's lecture notes, U. Michigan, 2014}.
\newblock [\epfmtdoi{10.1080/01621459.1963.10500830}]

\bibitem{LZ17}\bibhead{LZ17}
{\sc Chengyu Lin and Shengyu Zhang}: Sensitivity conjecture and log-rank
  conjecture for functions with small alternating numbers.
\newblock In {\em Proc. 44th Internat. Colloq. on Automata, Languages, and
  Programming (ICALP'17)}, pp. 51:1--13. Schloss Dagstuhl--Leibniz-Zentrum fuer
  Informatik, 2017.
\newblock [\epfmtdoi{10.4230/LIPIcs.ICALP.2017.51}]

\bibitem{LMN93}\bibhead{LMN93}
{\sc Nathan Linial, Yishay Mansour, and Noam Nisan}: Constant depth circuits,
  {F}ourier transform, and learnability.
\newblock {\em J. ACM}, 40(3):607--620, 1993.
\newblock [\epfmtdoi{10.1145/174130.174138}]

\bibitem{Mark}\bibhead{Mark}
{\sc Andrey~A Markov}: On the inversion complexity of a system of functions.
\newblock {\em J. ACM}, 5(4):331--334, 1958.
\newblock [\epfmtdoi{10.1145/320941.320945}]

\bibitem{book}\bibhead{book}
{\sc Ryan O'Donnell}: {\em Analysis of {B}oolean Functions}.
\newblock Cambridge Univ. Press, 2014.
\newblock [\epfmtdoi{10.1017/CBO9781139814782}, \epfmt{arxiv}{2105.10386}]

\bibitem{kkl09}\bibhead{kkl09}
{\sc Ryan O'Donnell and Karl Wimmer}: {KKL, Kruskal-Katona, and monotone nets}.
\newblock {\em SIAM J. Comput.}, 42(6):2375--2399, 2013.
\newblock \prelim \href{https://doi.org/10.1109/FOCS.2009.78}{FOCS'09}.
\newblock [\epfmtdoi{10.1137/100787325}]

\bibitem{razborov1985lower}\bibhead{razborov1985lower}
{\sc Alexander~A Razborov}: Lower bounds for the monotone complexity of some
  {B}oolean functions.
\newblock {\em Soviet Math. Doklady}, 31:354--357, 1985.
\newblock
  \href{https://www.mathnet.ru/links/c9ed090b71ff8f206b8e1a08ccc6cea4/dan9192.pdf}{Russian
  original at Mathnet.ru}.

\bibitem{Ros15}\bibhead{Ros15}
{\sc Benjamin Rossman}: Correlation bounds against monotone {NC$^1$}.
\newblock In {\em Proc. 30th Comput. Complexity Conf. (CCC'15)}, pp. 392--411.
  Schloss Dagstuhl--Leibniz-Zentrum fuer Informatik, 2015.
\newblock [\epfmtdoi{10.4230/LIPIcs.CCC.2015.392}]

\bibitem{srinivasan2011symmetric}\bibhead{srinivasan2011symmetric}
{\sc Murali~K Srinivasan}: {Symmetric chains, Gelfand--Tsetlin chains, and the
  Terwilliger algebra of the binary Hamming scheme}.
\newblock {\em J. Algebraic Combinatorics}, 34(2):301--322, 2011.
\newblock [\epfmtdoi{10.1007/s10801-010-0272-2}]

\bibitem{tardos1988gap}\bibhead{tardos1988gap}
{\sc {\'E}va Tardos}: The gap between monotone and non-monotone circuit
  complexity is exponential.
\newblock {\em Combinatorica}, 8(1):141--142, 1988.
\newblock [\epfmtdoi{10.1007/BF02122563}]

\end{thebibliography}

\end{document}
